% GENERATED FILE. Edit canonical lessons, then run npm run generate:books.
% canonical-source-hash: fnv1a64:ae5a706a043ba86a
% canonical-lessons: MR-C554-muddam, MR-C554-agau, MR-C554-ainveli, MR-C554-bahutek, MR-C554-sadharan

\chapter{Chance and Timing}
\label{ch:mr-chance-and-timing}
\hlchaptermodality{554}

\begin{chapteropening}
\textbf{By the end of this chapter:} \emph{I can say when and how something happens.}

Complete the last lesson of chapter 554: I can say when and how something happens.
\end{chapteropening}

\section[\texorpdfstring{\mr{मुद्दाम} (muddām)}{मुद्दाम (muddām)}]{\mr{मुद्दाम} (muddām) — on purpose}

\label{lesson:MR-C554-muddam}



\begin{quote}
Before the new one: say the Marathi for without fail, then the Marathi for prompt.
\end{quote}

\subsection*{You'll want to know: \mr{मुद्दाम}}
\textbf{\mr{मुद्दाम}} — \emph{muddām} — \textquotedblleft{}on purpose\textquotedblright{}.

\begin{quote}
\mr{मी} \mr{मुद्दाम} \mr{उशिरा} \mr{आलो} \mr{नाही}.
\end{quote}

\begin{quote}
\emph{mī muddām uśirā ālo nāhī.}
\end{quote}

\begin{quote}
I did not come late on purpose.
\end{quote}

\begin{cousinweb}
\textbf{\mr{मुद्दाम}} (\emph{muddām}) is a Marathi adverb for something done deliberately; its origin is not settled.
\end{cousinweb}

\begin{grammarlens}[title={What you've built}]
One more word for saying what happens if.
\end{grammarlens}

\subsection*{Your turn}
\begin{itemize}
\raggedright
  \item \emph{Say it:} \emph{muddām}
  \item \emph{Say it:} \emph{muddām}, once more
  \item \emph{Say it:} \emph{mī muddām uśirā ālo nāhī.} — I did not come late on purpose.
  \item \emph{Recall:} say the Marathi for without fail, then the Marathi for prompt, then say \emph{muddām} again
\end{itemize}

\begin{tcolorbox}[breakable,colback=teal!4,colframe=teal!35!black,title={Before you move on}]
What does \mr{मुद्दाम} mean? (\textquotedblleft{}On purpose\textquotedblright{}.) Say it once more.
\end{tcolorbox}
\section[\texorpdfstring{\mr{आगाऊ} (āgāū)}{आगाऊ (āgāū)}]{\mr{आगाऊ} (āgāū) — in advance}

\label{lesson:MR-C554-agau}



\begin{quote}
Before the new one: say the Marathi for prompt, then the Marathi for on purpose.
\end{quote}

\subsection*{You'll want to know: \mr{आगाऊ}}
\textbf{\mr{आगाऊ}} — \emph{āgāū} — \textquotedblleft{}in advance\textquotedblright{}.

\begin{quote}
\mr{सहलीचे} \mr{पैसे} \mr{आम्ही} \mr{आगाऊ} \mr{भरले}.
\end{quote}

\begin{quote}
\emph{sahlīce paise āmhī āgāū bharle.}
\end{quote}

\begin{quote}
We paid for the trip in advance.
\end{quote}

\begin{cousinweb}
\textbf{\mr{आगाऊ}} (\emph{āgāū}), \textquotedblleft{}ahead of time\textquotedblright{}, is perhaps kin to Marathi \emph{agāv}, \textquotedblleft{}in front\textquotedblright{}, which CDIAL sets under Sanskrit \emph{agra}, \textquotedblleft{}top, front\textquotedblright{} (CDIAL 68). CDIAL lists no \mr{आगाऊ} itself, so the link is not settled.
\end{cousinweb}

\begin{grammarlens}[title={What you've built}]
One more word for saying what happens if.
\end{grammarlens}

\subsection*{Your turn}
\begin{itemize}
\raggedright
  \item \emph{Say it:} \emph{āgāū}
  \item \emph{Say it:} \emph{āgāū}, once more
  \item \emph{Say it:} \emph{sahlīce paise āmhī āgāū bharle.} — We paid for the trip in advance.
  \item \emph{Recall:} say the Marathi for prompt, then the Marathi for on purpose, then say \emph{āgāū} again
\end{itemize}

\begin{tcolorbox}[breakable,colback=teal!4,colframe=teal!35!black,title={Before you move on}]
What does \mr{आगाऊ} mean? (\textquotedblleft{}In advance\textquotedblright{}.) Say it once more.
\end{tcolorbox}
\section[\texorpdfstring{\mr{ऐनवेळी} (ainveḷī)}{ऐनवेळी (ainveḷī)}]{\mr{ऐनवेळी} (ainveḷī) — at the last moment}

\label{lesson:MR-C554-ainveli}



\begin{quote}
Before the new one: say the Marathi for on purpose, then the Marathi for in advance.
\end{quote}

\subsection*{You'll want to know: \mr{ऐनवेळी}}
\textbf{\mr{ऐनवेळी}} — \emph{ainveḷī} — \textquotedblleft{}at the last moment\textquotedblright{}.

\begin{quote}
\mr{ऐनवेळी} \mr{गाडी} \mr{बिघडली} \mr{आणि} \mr{आम्ही} \mr{उशिरा} \mr{पोहोचलो}.
\end{quote}

\begin{quote}
\emph{ainveḷī gāḍī bighaḍlī āṇi āmhī uśirā pohoclo.}
\end{quote}

\begin{quote}
At the last moment the car broke down and we arrived late.
\end{quote}

\begin{cousinweb}
\textbf{\mr{ऐनवेळी}} (\emph{ainveḷī}) is a Marathi compound of \mr{ऐन}, \textquotedblleft{}the very\textquotedblright{}, a loan usually traced to Arabic \emph{ʿain}, \textquotedblleft{}eye, essence\textquotedblright{}, and \mr{वेळ}, \textquotedblleft{}time\textquotedblright{}.
\end{cousinweb}

\begin{grammarlens}[title={What you've built}]
One more word for saying what happens if.
\end{grammarlens}

\subsection*{Your turn}
\begin{itemize}
\raggedright
  \item \emph{Say it:} \emph{ainveḷī}
  \item \emph{Say it:} \emph{ainveḷī}, once more
  \item \emph{Say it:} \emph{ainveḷī gāḍī bighaḍlī āṇi āmhī uśirā pohoclo.} — At the last moment the car broke down and we arrived late.
  \item \emph{Recall:} say the Marathi for on purpose, then the Marathi for in advance, then say \emph{ainveḷī} again
\end{itemize}

\begin{tcolorbox}[breakable,colback=teal!4,colframe=teal!35!black,title={Before you move on}]
What does \mr{ऐनवेळी} mean? (\textquotedblleft{}At the last moment\textquotedblright{}.) Say it once more.
\end{tcolorbox}
\section[\texorpdfstring{\mr{बहुतेक} (bahutek)}{बहुतेक (bahutek)}]{\mr{बहुतेक} (bahutek) — mostly; most}

\label{lesson:MR-C554-bahutek}



\begin{quote}
Before the new one: say the Marathi for in advance, then the Marathi for at the last moment.
\end{quote}

\subsection*{You'll want to know: \mr{बहुतेक}}
\textbf{\mr{बहुतेक}} — \emph{bahutek} — \textquotedblleft{}mostly; most\textquotedblright{}.

\begin{quote}
\mr{बहुतेक} \mr{लोक} \mr{रविवारी} \mr{घरी} \mr{असतात}.
\end{quote}

\begin{quote}
\emph{bahutek lok ravivārī gharī astāt.}
\end{quote}

\begin{quote}
Most people are at home on Sundays.
\end{quote}

\begin{cousinweb}
\textbf{\mr{बहुतेक}} (\emph{bahutek}) is a Marathi compound built from \mr{बहुत}, \textquotedblleft{}much\textquotedblright{}, and \mr{एक}, \textquotedblleft{}one\textquotedblright{}. CDIAL lists Marathi \emph{bahut}, \textquotedblleft{}much, many\textquotedblright{}, under Sanskrit \emph{bahutva}, \textquotedblleft{}abundance\textquotedblright{} (CDIAL 9190).
\end{cousinweb}

\begin{grammarlens}[title={What you've built}]
One more word for saying what happens if.
\end{grammarlens}

\subsection*{Your turn}
\begin{itemize}
\raggedright
  \item \emph{Say it:} \emph{bahutek}
  \item \emph{Say it:} \emph{bahutek}, once more
  \item \emph{Say it:} \emph{bahutek lok ravivārī gharī astāt.} — Most people are at home on Sundays.
  \item \emph{Recall:} say the Marathi for in advance, then the Marathi for at the last moment, then say \emph{bahutek} again
\end{itemize}

\begin{tcolorbox}[breakable,colback=teal!4,colframe=teal!35!black,title={Before you move on}]
What does \mr{बहुतेक} mean? (\textquotedblleft{}Mostly; most\textquotedblright{}.) Say it once more.
\end{tcolorbox}
\section[\texorpdfstring{\mr{साधारण} (sādhāraṇ)}{साधारण (sādhāraṇ)}]{\mr{साधारण} (sādhāraṇ) — roughly; ordinary}

\label{lesson:MR-C554-sadharan}



\begin{quote}
Before the new one: say the Marathi for at the last moment, then the Marathi for mostly.
\end{quote}

\subsection*{You'll want to know: \mr{साधारण}}
\textbf{\mr{साधारण}} — \emph{sādhāraṇ} — \textquotedblleft{}roughly; ordinary\textquotedblright{}.

\begin{quote}
\mr{बस} \mr{साधारण} \mr{दहा} \mr{मिनिटांनी} \mr{येईल}.
\end{quote}

\begin{quote}
\emph{bas sādhāraṇ dahā miniṭānnī yeīl.}
\end{quote}

\begin{quote}
The bus will come in roughly ten minutes.
\end{quote}

\begin{cousinweb}
\textbf{\mr{साधारण}} (\emph{sādhāraṇ}) is a tatsama, taken from Sanskrit \emph{sādhāraṇa}, \textquotedblleft{}common to all, general\textquotedblright{} (Monier-Williams).
\end{cousinweb}

\begin{grammarlens}[title={What you've built}]
One more word for saying what happens if.
\end{grammarlens}

\subsection*{Your turn}
\begin{itemize}
\raggedright
  \item \emph{Say it:} \emph{sādhāraṇ}
  \item \emph{Say it:} \emph{sādhāraṇ}, once more
  \item \emph{Say it:} \emph{bas sādhāraṇ dahā miniṭānnī yeīl.} — The bus will come in roughly ten minutes.
  \item \emph{Recall:} say the Marathi for at the last moment, then the Marathi for mostly, then say \emph{sādhāraṇ} again
\end{itemize}

\begin{tcolorbox}[breakable,colback=teal!4,colframe=teal!35!black,title={Before you move on}]
What does \mr{साधारण} mean? (\textquotedblleft{}Roughly; ordinary\textquotedblright{}.) Say it once more.
\end{tcolorbox}
